\section{The Robustness--Retention Trade-off: A Formal Impossibility Result}
\label{sec:tradeoff}

We now formalize the coupling observed in~\cref{sec:motivation} as a provable lower bound on the trade-off between erasure robustness and neighbor preservation. The governing quantity is the \emph{concept overlap coefficient} $\kappa$, which measures the fraction of the target concept's activation region shared with its semantic neighbors.
% The empirical observations in Section~\ref{sec:motivation} suggest that concept leakage and collateral forgetting are coupled consequences of representational overlap. In this section, we make this intuition precise. We introduce a formal framework based on, what we define as, \emph{concept activation regions} in the model's internal representation space and prove that any unlearning method faces a fundamental lower bound on the trade-off between erasure robustness and neighbor preservation. The key quantity governing this trade-off is the \emph{concept overlap coefficient}~$\kappa$, which measures the fraction of the target concept's activation region that is shared with its semantic neighbors.

% \subsection{Setup and Definitions}
% \label{subsec:definitions}
% %\ali{the notation for the text encorder doesn't match the one used earlier in the background}
% Let $\mathcal{D}_\theta$ be a pretrained text-to-image diffusion model with U{-}Net noise predictor $\boldsymbol{\epsilon}_\theta$ and text encoder $f_\psi : \mathcal{P} \to \mathbb{R}^d$, where $\mathcal{P}$ is the space of text prompts. 
% For a prompt $p \in \mathcal{P}$, we denote its embedding by $\mathbf{e}_p = f_\psi(p) \in \mathbb{R}^d$. 
% And we write $p_\theta(\mathbf{x} \mid p)$ for the conditional distribution over images induced by $\mathcal{D}_\theta$.
% Let $\mathcal{D}_{\hat\theta}$ denote a concept-erased model (CEM) obtained by applying an unlearning procedure. %that modifies the U-Net parameters $\theta \mapsto \hat\theta$,
% %\varun{use tilde for all unlearning related things? don't switch b/w hat and tilde?}
% We correspondingly write $\boldsymbol{\epsilon}_{\hat\theta}$ for the edited U-Net and $p_{\hat\theta}(\mathbf{x} \mid p)$ for the induced conditional distribution. %\varun{part that bugs me is that $\theta$ is used for both U-Net and full diffusion model?}
%  % \ali{I think there are still slight notation inconsistency with the evaluation section. For example the unlearned model are shown with $\mathcal{M_u}$ there.}

% % We model a concept through the set of internal activations that arise when the model generates images containing that concept. Concretely, for a fixed UNet layer and denoising timestep, we extract the corresponding cross-attention activation induced by a text prompt. This lets us represent each concept as a region in the model's internal activation space.

% \paragraph{Concept activation regions.}
% For a fixed probe layer $\ell^*$ and denoising timestep $t^*$, let $\boldsymbol{\Phi}_\theta(p)$ denote the mean-pooled cross-attention activation produced by prompt $p \in \mathcal{P}$, where $\mathcal{H} := \mathbb{R}^{d_{\ell^*}}$ is the activation space at layer $\ell^*$ and $d_{\ell^*}$ is the per-token channel dimension at that layer. Each concept $c \in \mathcal{C}$ is associated with its \emph{activation region} $\mathcal{R}_c \subseteq \mathcal{H}$ produced by prompts whose generated images contain $c$. The formal construction is given in~\cref{app:concept_regions}.

% \noindent The key quantity governing the trade-off is the degree of representational overlap between the target concept and other concepts.
% To make this precise, for sets $A, B \subseteq \mathcal{H}$ let
% $d_H(A, B) \;=\; \max\!\left\{\,\sup_{\mathbf{a} \in A}\inf_{\mathbf{b} \in B}\|\mathbf{a}-\mathbf{b}\|_2,\;\; \sup_{\mathbf{b} \in B}\inf_{\mathbf{a} \in A}\|\mathbf{a}-\mathbf{b}\|_2\,\right\}$
% denote the Hausdorff distance between $A$ and $B$ in $\mathcal{H}$, and for $\rho > 0$ let
% \begin{equation}
%     \mathcal{R}_c^{(\rho)}
%     \;=\;
%     \bigl\{\mathbf{h} \in \mathcal{H} \;\big|\;
%     \inf_{\mathbf{r}  \in \mathcal{R}_c}\|\mathbf{h} - \mathbf{r}\|_2 \leq \rho\bigr\}
% \end{equation}
% denote the closed $\rho$-neighborhood of $\mathcal{R}_c$ in $\mathcal{H}$. 




%In practice, we estimate $d_H$ from finite activation samples using cosine distances on sample means, which serve as a stable, computationally tractable lower bound on Hausdorff distance under unit-norm activations (see~\cref{app:kappa_estimation}).

% ----- uncomment if have space -----
% \begin{figure}[t]
% \centering
% \begin{tikzpicture}[scale=1.2]
%   % Fattened (dashed) regions — these overlap to form O
%   \draw[dashed, thick, blue!60] (-1.2,0) ellipse (1.5 and 1.1);
%   \draw[dashed, thick, red!60]  (1.2,0)  ellipse (1.5 and 1.1);
%   % Overlap region O (intersection of the two dashed ellipses), shaded
%   \begin{scope}
%     \clip (-1.2,0) ellipse (1.5 and 1.1);
%     \fill[purple!30] (1.2,0) ellipse (1.5 and 1.1);
%   \end{scope}
%   % Core (solid) regions — DISJOINT, close but not overlapping
%   \draw[thick, blue, fill=blue!12] (-1.2,0) ellipse (1.1 and 0.75);
%   \draw[thick, red,  fill=red!12]  (1.2,0)  ellipse (1.1 and 0.75);
%   % Hausdorff distance — between rightmost point of R_{c_u} and leftmost of R_{c'}
%   \draw[<->, thick] (-0.1,-0.05) -- (0.1,-0.05);
%   \node[font=\footnotesize, above] at (0,-0.05) {$d_H$};
%   % rho annotation — vertical double-arrow on top of blue ellipse
%   \draw[<->, thick, gray] (-1.2,0.75) -- (-1.2,1.1);
%   \node[font=\footnotesize, gray, right] at (-1.2,0.92) {$\rho$};
%   % Displacement arrow Delta — from a point inside O outward into R_{c'}
%   \draw[->, very thick, purple!80!black] (0.05,0.25) -- (0.95,0.25);
%   \node[font=\footnotesize, purple!80!black, above] at (0.5,0.27) {$\Delta$};
%   % Region labels
%   \node[blue!70!black] at (-1.2,0) {$\mathcal{R}_{c_u}$};
%   \node[red!70!black]  at (1.2,0)  {$\mathcal{R}_{c'}$};
%   \node[font=\footnotesize, blue!60] at (-2.0,1.25) {$\mathcal{R}_{c_u}^{(\rho)}$};
%   \node[font=\footnotesize, red!60]  at (2.0,1.25)  {$\mathcal{R}_{c'}^{(\rho)}$};
%   \node[font=\footnotesize, purple!70!black] at (0,-0.55) {$\mathcal{O}$};
%   % kappa annotation in lower-left
%   \node[font=\footnotesize, align=left] at (-2.4,-1.5) 
%     {$\kappa = \dfrac{\mathrm{vol}(\mathcal{O})}{\mathrm{vol}(\mathcal{R}_{c_u}^{(\rho)})}$};
% \end{tikzpicture}
% \caption{Geometric setup of~\cref{thm:tradeoff}. Solid ellipses depict activation regions $\mathcal{R}_{c_u}$ (target) and $\mathcal{R}_{c'}$ (neighbor); dashed boundaries show their $\rho$-fattenings. The Hausdorff distance $d_H$ is the worst-case point-to-set separation between the unfattened regions. The shaded overlap $\mathcal{O} = \mathcal{R}_{c_u}^{(\rho)} \cap \mathcal{R}_{c'}^{(\rho)}$ arises because the fattened regions intersect even though the underlying activation regions do not. An $\varepsilon$-robust unlearning edit displaces target activations by at least $\Delta = (1-\varepsilon)/L$; for points $\mathbf{h}^* \in \mathcal{O}$, this displacement propagates into $\mathcal{R}_{c'}$, yielding the bound $\gamma \geq \kappa\Delta - d_H$.}
% \label{fig:geometry}
% \end{figure}

% , which preserves the set-theoretic meaning of the definition under unit normalization (see remark in~\cref{app:proof}).
% The fattening is a standard regularization that accounts for the fact that activation regions are estimated from finite prompt samples and may not be convex or closed.


% \begin{definition}[Concept entanglement coefficient]
% \label{def:overlap}
% The \emph{entanglement coefficient} of target concept $c_u$ with respect to its
% $\delta$-neighborhood and fattening radius $\rho > 0$ is
% \[
%     \kappa(c_u,\,\delta,\,\rho)
%     \;=\;
%     \frac{
%         \mathrm{vol}\!\left(
%             \mathcal{R}_{c_u}^{(\rho)}
%             \cap
%             \bigcup_{c' \in \mathcal{N}_\delta(c_u)} \mathcal{R}_{c'}^{(\rho)}
%             % \bigcup_{\substack{c' \in \mathcal{C},\; c' \neq c_u \\ d_{\mathcal{H}}(\mathcal{R}_{c_u},\,\mathcal{R}_{c'}) < \delta}}
% \mathcal{R}_{c'}^{(\rho)}
%         \right)
%     }{
%         \mathrm{vol}\!\left(\mathcal{R}_{c_u}^{(\rho)}\right)
%     },
% \]
% where $\mathrm{vol}(\cdot)$  denotes the $n$-dimensional volume of the $\rho$-fattened regions in $\mathcal{H} \cong \mathbb{R}^n$ and  $\mathcal{N}_\delta(c_u) = \{c' \in \mathcal{C} \mid c' \neq c_u,\; d_{\mathcal{H}} (\mathcal{R}_{c_u}, \mathcal{R}_{c'}) < \delta\}$  \ali{can we completely get rid of the neighborhood notion and just do a union over all the remaining concepts?}denotes the set of $\delta$-neighbors of $c_u$.
% We write $\kappa(c_u, \delta)$ when $\rho$ is held fixed.

% % \ali{why do we need a separate definition of neighborhood? we can directly do a union over all the concepts in this definition and get rid of Definition 4.2. less definitions makes it easier to follow.}
% \end{definition}



% \begin{definition}[Concept entanglement coefficient]
% \label{def:overlap}
% The \emph{entanglement coefficient} of target concept $c_u$ with respect to its
% $\delta$-neighborhood and fattening radius $\rho > 0$ is
% \[
%     \kappa(c_u,\,\delta,\,\rho)
%     \;=\;
%     \frac{
%         \mathrm{vol}\!\left(
%             \mathcal{R}_{c_u}^{(\rho)}
%             \cap
%             \bigcup_{c' \in \mathcal{C}} \mathcal{R}_{c'}^{(\rho)}
%             % \bigcup_{\substack{c' \in \mathcal{C},\; c' \neq c_u \\ d_{\mathcal{H}}(\mathcal{R}_{c_u},\,\mathcal{R}_{c'}) < \delta}}
% %\mathcal{R}_{c'}^{(\rho)}
%         \right)
%     }{
%         \mathrm{vol}\!\left(\mathcal{R}_{c_u}^{(\rho)}\right)
%     },
% \]
% where $\mathrm{vol}(\cdot)$  denotes the $n$-dimensional volume of the $\rho$-fattened regions in $\mathcal{H} \cong \mathbb{R}^n$ 
% % and  $\mathcal{N}_\delta(c_u) = \{c' \in \mathcal{C} \mid c' \neq c_u,\; d_{\mathcal{H}} (\mathcal{R}_{c_u}, \mathcal{R}_{c'}) < \delta\}$  \ali{can we completely get rid of the neighborhood notion and just do a union over all the remaining concepts?}
% %denotes the set of $\delta$-neighbors of $c_u$.
% We write $\kappa(c_u, \delta)$ when $\rho$ is held fixed.

% % \ali{why do we need a separate definition of neighborhood? we can directly do a union over all the concepts in this definition and get rid of Definition 4.2. less definitions makes it easier to follow.}
% \end{definition}
% \noindent Intuitively, $\kappa \in [0,1]$ measures what fraction of the target concept's representational footprint is shared with neighboring concepts. When $\kappa = 0$, the concept is fully disentangled from its neighbors and can in principle be erased without collateral damage. When $\kappa > 0$, any edit that displaces activations in $\mathcal{R}_{c_u}$ will inevitably disturb the overlapping portions of $\mathcal{R}_{c'}$.
% This formalizes the notion of entanglement introduced in Section 3 in terms of activation-space overlap.
% \varun{i love all of this empirical evidence; but please move to appendix and provide a single line with the main summary/takeaway and a pointer to the appendix!}



\subsection{Setup and Definitions}
\label{subsec:definitions}

Let $\mathcal{D}_\theta$ be a pretrained text-to-image diffusion model with U{-}Net noise predictor $\boldsymbol{\epsilon}_\theta$ and text encoder $f_\psi : \mathcal{P} \to \mathbb{R}^d$, where $\mathcal{P}$ is the space of text prompts. 
For a prompt $p \in \mathcal{P}$, we denote its embedding by $\mathbf{e}_p = f_\psi(p) \in \mathbb{R}^d$, and write $p_\theta(\mathbf{x} \mid p)$ for the conditional distribution over images induced by $\mathcal{D}_\theta$.
Let $\mathcal{D}_{\hat\theta}$ denote a concept-erased model (CEM) obtained by applying an unlearning procedure. 
We correspondingly write $\boldsymbol{\epsilon}_{\hat\theta}$ for the edited U-Net and $p_{\hat\theta}(\mathbf{x} \mid p)$ for the induced conditional distribution. 

\paragraph{Concept activation regions.}
For a fixed probe layer $\ell^*$ and denoising timestep $t^*$, let $\boldsymbol{\Phi}_\theta(p)$ denote the mean-pooled cross-attention activation produced by prompt $p \in \mathcal{P}$, where $\mathcal{H} := \mathbb{R}^{d_{\ell^*}}$ is the activation space at layer $\ell^*$ and $d_{\ell^*}$ is the per-token channel dimension at that layer. Each concept $c \in \mathcal{C}$ is associated with its \emph{activation region} $\mathcal{R}_c \subseteq \mathcal{H}$ produced by prompts whose generated images contain $c$. The formal construction is given in~\cref{app:concept_regions}.
Let
\begin{equation}
    \mathcal{R}_c^{(\rho)}
    \;=\;
    \bigl\{\mathbf{h} \in \mathcal{H} \;\big|\;
    \inf_{\mathbf{r}  \in \mathcal{R}_c}\|\mathbf{h} - \mathbf{r}\|_2 \leq \rho\bigr\}
\end{equation}
denote the closed $\rho$-neighborhood of $\mathcal{R}_c$ in $\mathcal{H}$. The key quantity governing the trade-off is the degree of representational overlap between the target concept and other concepts, defined as follows. 


\begin{definition}
\label{def:overlap}
The \emph{entanglement coefficient} of target concept $c_u$ with fattening radius $\rho > 0$ is
\[
    \kappa(c_u,\,\rho)
    \;=\;
    \frac{
        \mathrm{vol}\!\left(
            \mathcal{R}_{c_u}^{(\rho)}
            \cap
            \bigcup_{c' \in \mathcal{C} \setminus \{c_u\}} \mathcal{R}_{c'}^{(\rho)}
        \right)
    }{
        \mathrm{vol}\!\left(\mathcal{R}_{c_u}^{(\rho)}\right)
    },
\]
where $\mathrm{vol}(\cdot)$ denotes the $n$-dimensional volume of subsets of $\mathcal{H} \cong \mathbb{R}^n$.
\end{definition}

\noindent Intuitively, $\kappa \in [0,1]$ measures what fraction of the target concept's representational footprint overlaps with that of other concepts in the model. When $\kappa = 0$, the concept is fully disentangled and can in principle be erased without collateral damage. When $\kappa > 0$, any edit that displaces activations in $\mathcal{R}_{c_u}$ will inevitably disturb overlapping portions of other concepts' activation regions.
This formulation aggregates overlap across all concepts. In practice, however, the contribution to $\kappa$ is dominated by concepts that occupy nearby regions of the activation space; distant concepts contribute negligibly to the intersection. 


% \paragraph{Quantifying entanglement.}
% \Cref{tab:kappa} reports the estimated entanglement coefficient $\hat\kappa$ for five target concepts under a fixed fattening radius $\rho$ and a per-concept neighbor set $\mathcal{N}_\delta(c_u)$ discovered from data (\cref{app:kappa_estimation}).
% % \varun{don't understand this statement -- entanglement is calculated w.r.t what?}
% Each $\hat\kappa(c_u)$ is computed against that concept's own discovered neighbor set, listed in the rightmost column of~\cref{tab:kappa}, under a single shared $(\delta,\rho)$ across rows.
% The four animal subtypes (\texttt{dog}, \texttt{bear}, \texttt{horse}, \texttt{cat}) overlap substantially with fine-grained synonyms (e.g., \texttt{donkey}, \texttt{mare} for \texttt{horse}), while castle's architectural neighbors occupy more distinct regions despite semantic relatedness. 
% This $\hat\kappa$ heterogeneity is precisely the dose-response axis along which~\Cref{thm:tradeoff} predicts collateral damage should scale: erasing a high-$\hat\kappa$ concept should incur substantially larger neighbor-preservation error than erasing a low-$\hat\kappa$ one. We test this prediction in~\cref{sec:experiment}. The small number of targets here reflects the cost of the per-target erasure benchmark in~\cref{sec:experiment}; estimation details and sensitivity analyses are in~\cref{app:kappa_estimation}.


\paragraph{Quantifying entanglement.}
\Cref{tab:kappa} reports the estimated lower-bound for entanglement coefficient $\hat\kappa$ for five target concepts under a fixed fattening radius $\rho$ (\cref{app:kappa_estimation}). Each $\hat\kappa(c_u)$ is computed by measuring the fraction of activation samples from $c_u$ that lie within distance $\rho$ of any activation from other concepts. 
The four animal subtypes (\texttt{dog}, \texttt{bear}, \texttt{horse}, \texttt{cat}) exhibit substantial overlap with fine-grained variants (e.g., \texttt{donkey}, \texttt{mare} for \texttt{horse}), while \texttt{castle} exhibits significantly lower overlap despite having semantically related concepts. 
This heterogeneity in $\hat\kappa$ defines the regime in which~\Cref{thm:tradeoff} predicts collateral damage should scale: erasing a high-$\hat\kappa$ concept incurs substantially larger preservation error than erasing a low-$\hat\kappa$ one. We test this prediction in~\cref{sec:experiment}. Estimation details and sensitivity analyses are in~\cref{app:kappa_estimation}.

We now define the two desiderata that any unlearning method must satisfy:
\begin{definition}[$\varepsilon$-robust unlearning]
\label{def:robust}
The CEM $\mathcal{D}_{\hat\theta}$ achieves \emph{$\varepsilon$-robust unlearning} of concept $c_u$ if, for all prompts $p \in \mathcal{P}$ (including indirect prompts that do not
explicitly mention $c_u$),
\begin{equation}
    \Pr_{\mathbf{x} \sim p_{\hat\theta}(\cdot \mid p)}
\!\left[
\mathbb{O}_{\mathcal{X}}(\mathbf{x} ,\, c_u) = 1
\right]
\;<\; \varepsilon.
\end{equation}
where $\mathbb{O}_{\mathcal{X}} : \mathcal{X} \times \mathcal{C} \to \{0,1\}$ is the oracle concept detector introduced in~\cref{app:concept_regions}. The worst-case quantifier reflects an adversarial threat model; in~\cref{sec:experiment} we additionally report the corresponding average-case quantity UA.
\end{definition}

% \varun{should be $x \sim \bar{f}_{\boldsymbol{\boldsymbol{\Phi}}}(.|p)$?}


\begin{definition}[$\gamma$-concept preservation]
\label{def:neighbor}
Let $\mathcal{P}_{c'} \subseteq \mathcal{P}$ denote the set of prompts associated with concept $c'$. 
The CEM $\mathcal{D}_{\hat\theta}$ achieves \emph{$\gamma$-concept preservation} if, for all $c' \in \mathcal{C}$ with $c' \neq c_u$,
\begin{equation}
    \mathbb{E}_{p \sim \mathcal{P}_{c'}}\!\left[\bigl\|\boldsymbol{\Phi}_{\hat\theta}(p) - \boldsymbol{\Phi}_\theta(p)\bigr\|_2\right]
    \;<\; \gamma.
\end{equation}
\end{definition}

\noindent \Cref{def:robust} captures robustness to indirect prompts, while~\cref{def:neighbor} captures preservation of non-target concepts. Our main result below shows these two requirements are in tension.

% \begin{definition}[$\gamma$-neighbor preservation]
% \label{def:neighbor}
% Let $\mathcal{P}_{c'} \subseteq \mathcal{P}$ denote the set of prompts associated with concept $c'$.
% % The CEM $\mathcal{D}_{\hat\theta}$ achieves \emph{$\gamma$-neighbor preservation} with respect to $\mathcal{N}_\delta(c_u)$ if, 
% % for all $c' \in \mathcal{N}_\delta(c_u)$,
% The CEM $\mathcal{D}_{\hat\theta}$ achieves \emph{$\gamma$-neighbor preservation} at radius $\delta$ if, for all $c' \in \mathcal{C}$ with $c' \neq c_u$ and $d_{\mathcal{H}}(\mathcal{R}_{c_u}, \mathcal{R}_{c'}) < \delta$,
% \begin{equation}
%     \mathbb{E}_{p \sim \mathcal{P}_{c'}}\!\left[\bigl\|\boldsymbol{\Phi}_{\hat\theta}(p) - \boldsymbol{\Phi}_\theta(p)\bigr\|_2\right]
%     \;<\; \gamma.
% \end{equation}
% \end{definition}

% \noindent \Cref{def:robust} captures robustness to indirect prompts while~\cref{def:neighbor} captures preservation of neighboring concepts. Our main result below shows these two requirements are in tension.

\subsection{Main Result: The trade-off Lower Bound}
\label{subsec:main_result}
We require two mild regularity assumptions. Let $g_\theta: \mathcal{H} \to [0, 1]$ denote the composed map 
\begin{equation}
    g_\theta(h; c) \;:=\; \Pr_{\mathbf{x} \sim p_\theta(\cdot \mid \boldsymbol{\Phi}_\theta^{-1}(h))}\!\left[\mathbb{O}_{\mathcal{X}}(\mathbf{x}, c) = 1\right],
\end{equation}
i.e., the probability that concept $c$ appears in an image generated from an activation $h \in \mathcal{H}$. 
% remaining U-Net denoising steps (from layer $\ell^*$ onward) and the VAE decoder $V$. By construction, concept generation probability at prompt $p$ under model $\theta$ equals $g_\theta(\boldsymbol{\Phi}_\theta(p); c)$.
\begin{assumption}[Lipschitz continuity of generation]
\label{asm:lip}
For each concept $c \in \mathcal{C}$, 
$g_\theta(\cdot; c)$ is $L$-Lipschitz in the cross-attention activation at layer $\ell^*$, in the following sense: for any two activations $h, h' \in \mathcal{H}$ fed into the remainder of the generation pipeline 
%(UNet layers after $\ell^*$, VAE decoder $V$, and concept oracle),
\begin{equation}
    \bigl|
    \Pr[\mathbb{O}_{\mathcal{X}}(\mathbf{x}_h, c) = 1]
    -
    \Pr[\mathbb{O}_{\mathcal{X}}(\mathbf{x}_{h'}, c) = 1]
    \bigr|
    \;\leq\;
    L\,\|h - h'\|_2.
\end{equation}
where $x_h$ denotes an image generated from activation $h$. We further assume the unlearning edit modifies only layers at or before $\ell^*$, so the downstream pipeline is identical between $\mathcal{D}_\theta$ and $\mathcal{D}_{\hat\theta}$.
\end{assumption}
% \yian{empirical results: change the prompts a little bit and see if distance on embedding space changes (4.1 + 4.6)}

% \varun{i love all of this empirical evidence; but please move to appendix and provide a single line with the main summary/takeaway and a pointer to the appendix!}
\begin{assumption}[Minimum displacement for erasure]
\label{asm:edit}
Any unlearning edit $\theta \mapsto \hat\theta$ that achieves $\varepsilon$-robust unlearning of $c_u$ must displace activations within $\mathcal{R}_{c_u}$ by at least $\Delta > 0$:
\begin{equation}
    \forall\, p \in \mathcal{P} \;\text{ s.t. }\;
    \boldsymbol{\Phi}_{\theta}(p) \in \mathcal{R}_{c_u}:\quad
    \bigl\|\boldsymbol{\Phi}_{\hat\theta}(p) - \boldsymbol{\Phi}_{\theta}(p)\bigr\|_2
    \;\geq\; \Delta.
\end{equation}
\end{assumption}
\noindent We empirically verify~\cref{asm:lip} across natural prompt variations and~\cref{asm:edit} across seven unlearning methods erasing \texttt{horse}; all measured pairs satisfy Lipschitz with a small empirical constant, 
% and displacement propagates to semantic neighbors more than to controls. 
and displacement propagates to other concepts through shared activation structure.
Full details are in~\cref{app:lipschitz,app:displacement}.

% All five methods produce substantial displacement on target prompts ($\Delta_{\mathrm{target}} \geq 3.2$), confirming that $\varepsilon$-robust erasure requires non-trivial activation displacement as posited by~\cref{asm:edit}.
% Crucially, the displacement propagates to semantic neighbors: for every method, $\Delta_{\mathrm{neighbor}} > \Delta_{\mathrm{control}}$, and the gap is largest for STEREO ($\Delta_{\mathrm{neighbor}} - \Delta_{\mathrm{control}} = 3.26$), the method with the strongest erasure robustness. 
% This is consistent with Theorem~\ref{thm:tradeoff}: displacement within $\mathcal{R}_{c_u}$ propagates to overlapping neighbor regions in proportion to $\kappa$, causing greater collateral damage for high-$\kappa$ neighbors than for low-$\kappa$ controls.

% \varun{i love all of this empirical evidence; but please move to appendix and provide a single line with the main summary/takeaway and a pointer to the appendix!}
% \begin{theorem}[Robustness--Retention trade-off]
% \label{thm:tradeoff}
% Let Assumptions~\ref{asm:lip} and~\ref{asm:edit} hold with constants $L$ and $\Delta$.
% Let $\rho > 0$ be the fattening radius from Definition~\ref{def:overlap}.
% Suppose $\mathcal{D}_{\hat\theta}$ achieves $\varepsilon$-robust unlearning of $c_u$
% (Definition~\ref{def:robust}).
% % Then for any neighbor $c' \in \mathcal{N}_\delta(c_u)$ with
% % $\mathcal{R}_{c'}^{(\rho)} \cap \mathcal{R}_{c_u}^{(\rho)} \neq \emptyset$,
% Then for any $c' \in \mathcal{C}$ with $c' \neq c_u$,
% $d_{\mathcal{H}}(\mathcal{R}_{c_u}, \mathcal{R}_{c'}) < \delta$, and
% $\mathcal{R}_{c'}^{(\rho)} \cap \mathcal{R}_{c_u}^{(\rho)} \neq \emptyset$,
% the neighbor preservation error satisfies
% \begin{equation}
%     \mathbb{E}_{p \sim \mathcal{P}_{c'}}\!\left[\bigl\|\boldsymbol{\Phi}_{\hat\theta}(p) - \boldsymbol{\Phi}_\theta(p)\bigr\|_2\right]
%     \;\geq\;
%     \kappa(c_u,\delta) \cdot \Delta
%     \;-\; \delta,
% \end{equation}
% and consequently, $\gamma$-neighbor preservation (Definition~\ref{def:neighbor}) requires
% \begin{equation}
% \label{eq:tradeoff}
%     \boxed{
%     \gamma \;\geq\; \kappa(c_u,\delta)\cdot\frac{1-\varepsilon}{L}
%     \;-\; \delta.
%     }
% \end{equation}
% Equivalently, achieving both $\varepsilon$-robustness and $\gamma$-neighbor preservation simultaneously requires
% \begin{equation}
% \label{eq:tradeoff_eps}
%     \varepsilon \;\geq\; 1 - \frac{L(\gamma + \delta)}{\kappa(c_u,\delta)}.
% \end{equation}
% \end{theorem}


\begin{theorem}[Robustness--Retention trade-off]
\label{thm:tradeoff}
Let Assumptions~\ref{asm:lip} and~\ref{asm:edit} hold with constants $L$ and $\Delta$.
Let $\rho > 0$ be the fattening radius from Definition~\ref{def:overlap}.
Suppose $\mathcal{D}_{\hat\theta}$ achieves $\varepsilon$-robust unlearning of $c_u$
(Definition~\ref{def:robust}).
%Then for any concept $c' \in \mathcal{C}$ with $c' \neq c_u$, the expected activation displacement satisfies
Then the aggregate collateral displacement over non-target concepts satisfies
\begin{equation}
    \mathbb{E}_{p\sim \mathcal{P}_{-u}}
    \left[
\|\boldsymbol{\Phi}_{\hat{\theta}}(p)-\boldsymbol{\Phi}_{\theta}(p)\|_2
    \right]
    \geq
    \kappa(c_u,\rho)\cdot \Delta .
\end{equation}
% where $\kappa(c_u,\rho)$ is the entanglement coefficient from Definition~\ref{def:overlap}.

where $\mathcal{P}_{-u}$ denote the distribution over non-target prompts. Consequently, achieving $\gamma$-concept preservation (Definition~\ref{def:neighbor}) requires
\begin{equation}
\label{eq:tradeoff}
    \boxed{
    \gamma \;\geq\; \kappa(c_u,\rho)\cdot\frac{1-\varepsilon}{L}.
    }
\end{equation}

Equivalently, achieving both $\varepsilon$-robustness and $\gamma$-concept preservation simultaneously requires
\begin{equation}
\label{eq:tradeoff_eps}
    \varepsilon \;\geq\; 1 - \frac{L\gamma}{\kappa(c_u,\rho)}.
\end{equation}
\end{theorem}





\begin{proof}[Proof sketch]
By~\cref{asm:lip}, %$\varepsilon$-robust unlearning requires displacing all activations in $\mathcal{R}_{c_u}$ by at least $\Delta = (1-\varepsilon)/L$. This displacement propagates to the fraction $\kappa(c_u,\delta)$ of $\mathcal{R}_{c'}$ that overlaps with $\mathcal{R}_{c_u}^{(\rho)}$, minus at most $\delta$ due to inter-concept distance. 
$\varepsilon$-robust unlearning requires displacing all activations in $\mathcal{R}_{c_u}$ by at least $\Delta = (1-\varepsilon)/L$. By Definition~\ref{def:overlap}, a fraction $\kappa(c_u,\rho)$ of the fattened region $\mathcal{R}_{c_u}^{(\rho)}$ overlaps with non-target concept regions. This displacement therefore affects at least a $\kappa(c_u,\rho)$ fraction of the shared activation mass, yielding an aggregate collateral displacement of at least $\kappa(c_u,\rho)\cdot \Delta$.
% For any neighbor $c'$ whose fattened activation region overlaps with $\mathcal{R}_{c_u}^{(\rho)}$, the displacement in the overlap region propagates to $\mathcal{R}_{c'}$, yielding a lower bound proportional to the overlap fraction $\kappa(c_u,\delta)$ minus the inter-concept distance $\delta$. Then substituting $\Delta = (1-\varepsilon)/L$ yields the stated bound on $\gamma$. 
Full proof is given in Appendix~\ref{app:proof}.
\end{proof}

\subsection{Consequences and Interpretation}
\label{subsec:consequences}

\paragraph{Impossibility of perfect unlearning with perfect retention.}
Based on~\Cref{thm:tradeoff}, simultaneous perfect erasure and perfect retention is infeasible whenever concepts are entangled:
\begin{corollary}[Impossibility result]
\label{cor:impossibility}
%If $\kappa(c_u, \delta) > 0$ and $\delta < \kappa(c_u,\delta)/L$, 
If $\kappa(c_u,\rho) > 0$,
then no unlearning method can simultaneously achieve $\varepsilon = 0$ (perfect erasure) and $\gamma = 0$ (perfect neighbor preservation).
\end{corollary}

\begin{proof}
Setting $\varepsilon = 0$ in Eq.~\eqref{eq:tradeoff} gives 
% $\gamma \geq \kappa(c_u,\delta)/L - \delta > 0$ 
$\gamma \geq \kappa(c_u,\rho)/L > 0$
under the stated conditions.
\end{proof}

\paragraph{The Pareto frontier.}
Eq.~\eqref{eq:tradeoff_eps} defines a \emph{feasibility boundary} in the $(\varepsilon, \gamma)$ plane: no method can simultaneously achieve an erasure robustness below $\varepsilon$ and a neighbor preservation error below $\gamma$ unless the pair $(\varepsilon, \gamma)$ lies above the curve
$\gamma = \kappa \cdot \frac{1-\varepsilon}{L}.$
Different unlearning methods correspond to different operating points on or above this curve: aggressive methods achieve small $\varepsilon$ at the cost of large $\gamma$, occupying the upper-left of the frontier; conservative methods accept large $\varepsilon$ to maintain small $\gamma$, occupying the lower-right.
In Section~\ref{sec:experiment}, we verify that the empirically observed operating points across seven methods are consistent with this predicted boundary.



% \paragraph{Empirical validation.}
% To verify that concept activation regions are well-defined empirical objects, we extract cross-attention activations from Stable Diffusion v1.4 at the mid-block layer (\texttt{mid\_block.attentions.0}) at step $t^*{=}25$ of 50 DDIM steps for 10 concepts, each prompted with 50 shared templates differing only in the concept word.
% Figure~\ref{fig:umap} shows that concepts with nearby activation centroids in the extracted cross-attention space from neighboring clusters. For example,: \texttt{horse}, \texttt{pony}, and \texttt{donkey} form a dense equine cluster, while \texttt{castle} is well-separated from all animal concepts. This provides empirical support that concept activation regions are coherent geometric objects in $\mathcal{H}$.


% Empirical evidence that these activation regions correspond to coherent geometric structure is shown in~\cref{fig:umap}; full construction details and discussion are deferred to Appendix~\ref{app:activation_regions}.

% \varun{this is EXCELLENT!!! i would move this figure to the intro; move the discussion of the figure (this empirical validation paragraph) to the appendix, but refer to the figure in a single line in this part of the paper to show that "empirical validation shows that these definitions are real"}


% We formalize the notion that concepts occupy regions in the model's internal activation space, and that these regions can overlap for semantically related concepts.

% \varun{what does this activation region capture? similarly what does the activation function capture?}
% \begin{definition}[Concept activation region]
% \label{def:activation}
% Fix a UNet layer $\ell^*$ and denoising timestep $t^* \in \{1,\ldots,T\}$. The \emph{concept activation function} is
% \[
%     \boldsymbol{\Phi}_\theta : \mathcal{T} \longrightarrow \mathcal{Z}, \qquad
%     \boldsymbol{\Phi}_\theta(p) \;=\; \mathrm{A}_\theta(p,\, \ell^*,\, t^*),
% \]
% where $\mathrm{A}_\theta(p,\ell,t)$ denotes the cross-attention output at layer $\ell$ and timestep $t$ \varun{what is $p$?}. The \emph{activation region} of concept $c$ is
% \[
%     \mathcal{R}_c \;=\; \bigl\{\, \boldsymbol{\Phi}_{\theta^*}(p) \;\big|\;
%     p \in \mathcal{T},\; \mathbb{O}_\mathcal{X}(f_\boldsymbol{\Phi}(p), c) = 1
%     \,\bigr\} \;\subseteq\; \mathcal{Z},
% \]
% where $\mathbb{O}_\mathcal{X} : \mathcal{X} \times \mathcal{C} \to \{0,1\}$ is the ground-truth oracle indicating whether concept $c$ is present in image $x$.
% \end{definition}

% \begin{definition}[Semantic neighborhood]
% \label{def:neighborhood}
% The \emph{semantic neighborhood} of concept $c$ at radius $\delta > 0$ is
% \[
%     \mathcal{N}_\delta(c) \;=\; \bigl\{\, c' \in \mathcal{C}
%     \;\big|\; c' \neq c, \;
%     d_\mathcal{Z}\!\left(\mathcal{R}_c,\, \mathcal{R}_{c'}\right)
%     < \delta \,\bigr\},
% \]
% where $d_\mathcal{Z}(A,B) = \inf_{a \in A, b \in B}\|a - b\|_2$ is the set distance in $\mathcal{Z}$. \varun{$\mathcal{R}_{c'}$ is not defined?}
% \end{definition}
% \begin{definition}[Semantic neighborhood]
% \label{def:neighborhood}
% For a concept $c \in \mathcal{C}$ and radius $\delta > 0$, its \emph{semantic neighborhood}
% is
% \begin{equation}
%     \mathcal{N}_{\delta}(c)
% \;=\;
% \left\{
% c' \in \mathcal{C}
% \;\middle|\;
% c' \neq c,\;
% d_{\mathcal{H}}(\mathcal{R}_c,\,\mathcal{R}_{c'}) < \delta
% \right\},
% \end{equation}
% where $\mathcal{R}_{c'}$ is the activation region of concept $c'$ as in
% Definition~\ref{def:activation}, and
% \begin{equation}
%     d_{\mathcal{H}}(A,B) \;=\; \inf_{a \in A,\; b \in B}\|a - b\|_2
% \end{equation}

% \ali{since we use infimum in the definition, Figure 1 might be misleading.}
% is the point set distance in $\mathcal{H}$.
% \end{definition}

% Figure~\ref{fig:distance_matrix} reports the pairwise centroid cosine distance between 10 example concept activation regions, ordered by hierarchical clustering to reveal block structure.
% Two patterns emerge that directly bear on the trade-off.
% First, concepts within the same semantic category have small centroid distances: \texttt{horse}$\leftrightarrow$\texttt{pony} ($0.158$), \texttt{horse}$\leftrightarrow$\texttt{donkey} ($0.329$), \texttt{dog}$\leftrightarrow$\texttt{cat} ($0.282$), \texttt{car}$\leftrightarrow$\texttt{truck} ($0.400$). These small distances suggest substantial representational overlap, predicting that erasing any concept in these clusters will displace activations into neighboring concepts' regions. Second, semantically isolated concepts such as \texttt{castle} have centroid distances exceeding $0.65$ to all animal concepts, implying $\kappa \approx 0$ and predicting minimal collateral damage under erasure. We verify these predictions empirically in~\cref{sec:experiment}.
% \ali{this seems like a good empirical evidence to me, but a reviewer might raise the concern than your definitions are based on infimum while the empricial distance is based on average values.}
% The four animal concepts cluster at high entanglement ($\hat\kappa \in [0.77, 0.89]$): their activation regions overlap substantially with fine-grained subtypes and near-synonyms in the same taxonomic category, such as \texttt{puppy} and \texttt{labrador} for \texttt{dog}, or \texttt{donkey}, \texttt{foal}, and \texttt{mare} for \texttt{horse}. \texttt{castle}, by contrast, exhibits substantially lower entanglement ($\hat\kappa = 0.27$) despite having a similar number of discovered architectural neighbors (\texttt{fortress}, \texttt{palace}, \texttt{citadel}, \texttt{keep}, \texttt{manor}): these concepts are semantically related but occupy more distinct regions of activation space. 
% \varun{next sentence comes out of the blue; perhaps do not need it here?}
% \Cref{thm:tradeoff} predicts that collateral damage from erasing a high-$\hat\kappa$ concept will be substantially greater than from erasing a low-$\hat\kappa$ one; we verify this in~\cref{sec:experiment}. Full estimation details are in~\cref{app:kappa_estimation}.
% \begin{table}[t]
% \centering
% \small
% \begin{tabular}{lccc}
% \toprule
% Concept & $\hat\kappa$ & Silhouette & Neighbors \\
% \midrule
% % pony    & 0.46 & 0.47 & horse, donkey \\
% dog     & 0.89 & 0.54 & puppy, labrador \\
% % donkey  & 0.04 & 0.66 & horse, pony, deer \\
% % car     & 0.02 & 0.59 & dog, cat, truck, tower \\
% % deer    & 0.02 & 0.71 & donkey \\
% bear    & 0.82 & 0.70 & black bear, brown bear, grizzly\\
% horse   & 0.81 & 0.43 & donkey,foal,mare \\
% cat     & 0.77 & 0.53 & kitten, panther \\
% % truck   & 0.00 & 0.76 & car, tower \\
% castle  & 0.27 & 0.89 & citadel, fortress, keep, manor, palace \\
% \bottomrule
% \end{tabular}
% \caption{Estimated entanglement coefficient $\hat\kappa$ and  silhouette score \varun{what is this?} for each concept, computed on mean-pooled 
% cross-attention activations with cosine distance. Higher $\hat\kappa$ indicates greater representational overlap with semantic neighbors; higher silhouette 
% indicates a more isolated activation region. See~\cref{app:kappa_estimation} for estimation details.}
% \label{tab:kappa}
% \end{table}

